\documentclass[10pt,a4paper]{amsart}
\usepackage{amsmath,amssymb,amsthm,graphicx}
\usepackage{verbatim}
\newenvironment{reference}{\comment}{\endcomment}
%\newenvironment{reference}{}{}
\newenvironment{slogan}{\comment}{\endcomment}
\newenvironment{history}{\comment}{\endcomment}
% For commutative diagrams we use Xy-pic
\usepackage[all]{xy}

% We use 2cell for 2-commutative diagrams.
\xyoption{2cell}
% Theorem environments.
%
\theoremstyle{plain}
\newtheorem{theorem}[subsection]{Theorem}
\newtheorem{proposition}[subsection]{Proposition}
\newtheorem{lemma}[subsection]{Lemma}

\theoremstyle{definition}
\newtheorem{definition}[subsection]{Definition}
\newtheorem{example}[subsection]{Example}
\newtheorem{exercise}[subsection]{Exercise}
\newtheorem{situation}[subsection]{Situation}

\theoremstyle{remark}
\newtheorem{remark}[subsection]{Remark}
\newtheorem{remarks}[subsection]{Remarks}

\numberwithin{equation}{subsection}

% Macros
%
\def\lim{\mathop{\mathrm{lim}}\nolimits}
\def\colim{\mathop{\mathrm{colim}}\nolimits}
\def\Spec{\mathop{\mathrm{Spec}}}
\def\Hom{\mathop{\mathrm{Hom}}\nolimits}
\def\Ext{\mathop{\mathrm{Ext}}\nolimits}
\def\SheafHom{\mathop{\mathcal{H}\!\mathit{om}}\nolimits}
\def\SheafExt{\mathop{\mathcal{E}\!\mathit{xt}}\nolimits}
\def\Sch{\mathit{Sch}}
\def\Mor{\mathop{\mathrm{Mor}}\nolimits}
\def\Ob{\mathop{\mathrm{Ob}}\nolimits}
\def\Sh{\mathop{\mathit{Sh}}\nolimits}
\def\NL{\mathop{N\!L}\nolimits}
\def\CH{\mathop{\mathrm{CH}}\nolimits}
\def\proetale{{pro\text{-}\acute{e}tale}}
\def\etale{{\acute{e}tale}}
\def\QCoh{\mathit{QCoh}}
\def\Ker{\mathop{\mathrm{Ker}}}
\def\Im{\mathop{\mathrm{Im}}}
\def\Coker{\mathop{\mathrm{Coker}}}
\def\Coim{\mathop{\mathrm{Coim}}}

% Boxtimes
%
\DeclareMathSymbol{\boxtimes}{\mathbin}{AMSa}{"02}

%
% Macros for moduli stacks/spaces
%
\def\QCohstack{\mathcal{QC}\!\mathit{oh}}
\def\Cohstack{\mathcal{C}\!\mathit{oh}}
\def\Spacesstack{\mathcal{S}\!\mathit{paces}}
\def\Quotfunctor{\mathrm{Quot}}
\def\Hilbfunctor{\mathrm{Hilb}}
\def\Curvesstack{\mathcal{C}\!\mathit{urves}}
\def\Polarizedstack{\mathcal{P}\!\mathit{olarized}}
\def\Complexesstack{\mathcal{C}\!\mathit{omplexes}}
% \Pic is the operator that assigns to X its picard group, usage \Pic(X)
% \Picardstack_{X/B} denotes the Picard stack of X over B
% \Picardfunctor_{X/B} denotes the Picard functor of X over B
\def\Pic{\mathop{\mathrm{Pic}}\nolimits}
\def\Picardstack{\mathcal{P}\!\mathit{ic}}
\def\Picardfunctor{\mathrm{Pic}}
\def\Deformationcategory{\mathcal{D}\!\mathit{ef}}
\usepackage{cleveref}
\crefname{originallemma}{original lemma}{original lemmas}
\Crefname{originallemma}{Original lemma}{Original lemmas}
\makeatletter
\newlabel{derived-lemma-third-isomorphism-triangle}{{014A}{source}{}{sourcepointer-1}{}}
\newlabel{derived-lemma-third-isomorphism-triangle@cref}{{[originallemma][0][]014A}{[1][0][]source}{}{}{}}
\makeatother
\crefname{originalequation}{original equation}{original equations}
\Crefname{originalequation}{Original equation}{Original equations}
\makeatletter
\newlabel{equation-wedge}{{0FM2}{source}{}{sourcepointer-2}{}}
\newlabel{equation-wedge@cref}{{[originalequation][0][]0FM2}{[1][0][]source}{}{}{}}
\makeatother
\crefname{originallemma}{original lemma}{original lemmas}
\Crefname{originallemma}{Original lemma}{Original lemmas}
\makeatletter
\newlabel{lemma-duality-hodge}{{0FW4}{source}{}{sourcepointer-3}{}}
\newlabel{lemma-duality-hodge@cref}{{[originallemma][0][]0FW4}{[1][0][]source}{}{}{}}
\makeatother
\crefname{originallemma}{original lemma}{original lemmas}
\Crefname{originallemma}{Original lemma}{Original lemmas}
\makeatletter
\newlabel{lemma-relative-bottom-part-degenerates}{{0G8H}{source}{}{sourcepointer-4}{}}
\newlabel{lemma-relative-bottom-part-degenerates@cref}{{[originallemma][0][]0G8H}{[1][0][]source}{}{}{}}
\makeatother
\crefname{originallemma}{original lemma}{original lemmas}
\Crefname{originallemma}{Original lemma}{Original lemmas}
\makeatletter
\newlabel{lemma-relative-duality-hodge}{{0G8I}{source}{}{sourcepointer-5}{}}
\newlabel{lemma-relative-duality-hodge@cref}{{[originallemma][0][]0G8I}{[1][0][]source}{}{}{}}
\makeatother
\crefname{originallemma}{original lemma}{original lemmas}
\Crefname{originallemma}{Original lemma}{Original lemmas}
\makeatletter
\newlabel{lemma-relative-top-part-degenerates}{{0G8J}{source}{}{sourcepointer-6}{}}
\newlabel{lemma-relative-top-part-degenerates@cref}{{[originallemma][0][]0G8J}{[1][0][]source}{}{}{}}
\makeatother
\crefname{originallemma}{original lemma}{original lemmas}
\Crefname{originallemma}{Original lemma}{Original lemmas}
\makeatletter
\newlabel{limits-lemma-higher-direct-images-zero-above-dimension-fibre}{{0E7D}{source}{}{sourcepointer-7}{}}
\newlabel{limits-lemma-higher-direct-images-zero-above-dimension-fibre@cref}{{[originallemma][0][]0E7D}{[1][0][]source}{}{}{}}
\makeatother
\crefname{originalproposition}{original proposition}{original propositions}
\Crefname{originalproposition}{Original proposition}{Original propositions}
\makeatletter
\newlabel{proposition-poincare-duality}{{0FW7}{source}{}{sourcepointer-8}{}}
\newlabel{proposition-poincare-duality@cref}{{[originalproposition][0][]0FW7}{[1][0][]source}{}{}{}}
\makeatother
\crefname{originalsection}{original section}{original sections}
\Crefname{originalsection}{Original section}{Original sections}
\makeatletter
\newlabel{section-cup-product}{{0FM1}{source}{}{sourcepointer-9}{}}
\newlabel{section-cup-product@cref}{{[originalsection][0][]0FM1}{[1][0][]source}{}{}{}}
\makeatother
\crefname{originalsection}{original section}{original sections}
\Crefname{originalsection}{Original section}{Original sections}
\makeatletter
\newlabel{section-hodge-to-de-rham}{{0FM6}{source}{}{sourcepointer-10}{}}
\newlabel{section-hodge-to-de-rham@cref}{{[originalsection][0][]0FM6}{[1][0][]source}{}{}{}}
\makeatother
\title{Relative Poincare duality for the de Rham complex of a proper smooth morphism with nonempty equidimensional fibres}
\author{The Stacks Project Authors}
\date{stacks-project:derham.tex:local-snapshot-20261012; GNU FDL 1.2 or later}
\begin{document}
\maketitle
\noindent\textit{Editorial note.} All delivered mathematical text is the author's own, copied from the pinned Stacks Project chapter file derham.tex; no statement or proof step was rewritten, reordered or removed.\par\medskip
\noindent\textit{Editorial note.} Preamble provenance: the pinned chapter file opens with \textbackslash{}input\{preamble\} and carries no preamble text of its own, so the project's single Stacks-wide preamble staged at sources/stacks-derham/preamble.tex is excerpted verbatim as three spans, lines 8-12 (verbatim package and the project's own reference/slogan/history display environments), lines 14-18 (Xy-pic setup) and lines 154-217 (theorem declarations and mathematical macros). The wrapper class and the amsmath/amssymb/amsthm/graphicx base packages are supplied by the assembler. No macro was invented or redefined.\par\medskip
\noindent\textit{Editorial note.} Same-chapter pointers are located at their \textbackslash{}label lines inside the pinned chapter file; cross-chapter pointers are located in the project's own tag list staged in the snapshot (sources/stacks-chow/tags.txt), at the verbatim lines carrying their canonical tags.\par\medskip
\noindent\textit{Allowed standard prerequisites (editorial).}\begin{itemize}
\item The relative Hodge-to-de Rham spectral sequence. Same chapter derham.tex, Section \textbackslash{}ref\{section-hodge-to-de-rham\} (invoked at proof line 6348). It relates the relative de Rham complex to the sheaves of relative differentials, which is what identifies the extreme relative cohomology sheaves.
\item Vanishing of the higher direct images above the fibre dimension, and degeneration of the extreme parts of the relative spectral sequence. Limits, label limits-lemma-higher-direct-images-zero-above-dimension-fibre (invoked at proof line 6350) and Same chapter derham.tex, labels lemma-relative-bottom-part-degenerates and lemma-relative-top-part-degenerates (invoked at proof lines 6351 and 6352). Together with the vanishing of \$\textbackslash{}Omega\textasciicircum{}i\_\{X/S\}\$ for \$i > n\$ they give \$R\textasciicircum{}0f\_*\textbackslash{}Omega\textasciicircum{}\textbackslash{}bullet\_\{X/S\} = R\textasciicircum{}0f\_*\textbackslash{}mathcal\{O\}\_X\$ and \$R\textasciicircum{}nf\_*\textbackslash{}Omega\textasciicircum{}n\_\{X/S\} = R\textasciicircum{}\{2n\}f\_*\textbackslash{}Omega\textasciicircum{}\textbackslash{}bullet\_\{X/S\}\$.
\item The relative Hodge duality pairing in the extreme case. Same chapter derham.tex, label lemma-relative-duality-hodge (invoked at proof lines 6364 and 6412). It supplies the trace map whose base case gives perfect pairings between \$Rf\_*\textbackslash{}mathcal\{O\}\_X\$ and \$Rf\_*\textbackslash{}Omega\textasciicircum{}n[n]\$.
\item The wedge product on the total complex and the construction of relative cup products. Same chapter derham.tex, equation equation-wedge (invoked at proof line 6368) and Section \textbackslash{}ref\{section-cup-product\} (invoked at proof line 6388). The wedge map restricts to the maps \$\textbackslash{}gamma\_p\$ from the truncated total complex, whose induced products are the relative pairings studied in the induction.
\item Two-out-of-three for perfect pairings of complexes, in the form of the third isomorphism of a triangle. Derived Categories, label derived-lemma-third-isomorphism-triangle (invoked at proof line 6464). It replaces the five-lemma step of the absolute case: once the pairings on the outer rows are known perfect and the four pairs of maps are compatible, the pairing in the middle is perfect.
\end{itemize}\medskip
\ifdefined\hypertarget\hypertarget{sourcepointer-1}{}\fi
\noindent\textit{Source reference.} Original lemma 014A: sources/stacks-chow/tags.txt, source lines 1465--1465. Navigation pointer for \textbackslash{}ref\{derived-lemma-third-isomorphism-triangle\} used inside the counted span (line 6464); the argument that consumes it is retained in the span.\par\medskip
\ifdefined\hypertarget\hypertarget{sourcepointer-2}{}\fi
\noindent\textit{Source reference.} Original equation 0FM2: sources/stacks-derham/pinned-original-source.tex, source lines 495--495. Navigation pointer for \textbackslash{}ref\{equation-wedge\} used inside the counted span (line 6368); the argument that consumes it is retained in the span.\par\medskip
\ifdefined\hypertarget\hypertarget{sourcepointer-3}{}\fi
\noindent\textit{Source reference.} Original lemma 0FW4: sources/stacks-derham/pinned-original-source.tex, source lines 4496--4496. Navigation pointer for \textbackslash{}ref\{lemma-duality-hodge\} used inside the counted span (line 6460); the argument that consumes it is retained in the span.\par\medskip
\ifdefined\hypertarget\hypertarget{sourcepointer-4}{}\fi
\noindent\textit{Source reference.} Original lemma 0G8H: sources/stacks-derham/pinned-original-source.tex, source lines 5855--5855. Navigation pointer for \textbackslash{}ref\{lemma-relative-bottom-part-degenerates\} used inside the counted span (line 6351); the argument that consumes it is retained in the span.\par\medskip
\ifdefined\hypertarget\hypertarget{sourcepointer-5}{}\fi
\noindent\textit{Source reference.} Original lemma 0G8I: sources/stacks-derham/pinned-original-source.tex, source lines 5877--5877. Navigation pointer for \textbackslash{}ref\{lemma-relative-duality-hodge\} used inside the counted span (line 6364); the argument that consumes it is retained in the span.\par\medskip
\ifdefined\hypertarget\hypertarget{sourcepointer-6}{}\fi
\noindent\textit{Source reference.} Original lemma 0G8J: sources/stacks-derham/pinned-original-source.tex, source lines 5979--5979. Navigation pointer for \textbackslash{}ref\{lemma-relative-top-part-degenerates\} used inside the counted span (line 6352); the argument that consumes it is retained in the span.\par\medskip
\ifdefined\hypertarget\hypertarget{sourcepointer-7}{}\fi
\noindent\textit{Source reference.} Original lemma 0E7D: sources/stacks-chow/tags.txt, source lines 17778--17778. Navigation pointer for \textbackslash{}ref\{limits-lemma-higher-direct-images-zero-above-dimension-fibre\} used inside the counted span (line 6350); the argument that consumes it is retained in the span.\par\medskip
\ifdefined\hypertarget\hypertarget{sourcepointer-8}{}\fi
\noindent\textit{Source reference.} Original proposition 0FW7: sources/stacks-derham/pinned-original-source.tex, source lines 4690--4690. Navigation pointer for \textbackslash{}ref\{proposition-poincare-duality\} used inside the counted span (line 6341); the argument that consumes it is retained in the span.\par\medskip
\ifdefined\hypertarget\hypertarget{sourcepointer-9}{}\fi
\noindent\textit{Source reference.} Original section 0FM1: sources/stacks-derham/pinned-original-source.tex, source lines 483--483. Navigation pointer for \textbackslash{}ref\{section-cup-product\} used inside the counted span (line 6388); the argument that consumes it is retained in the span.\par\medskip
\ifdefined\hypertarget\hypertarget{sourcepointer-10}{}\fi
\noindent\textit{Source reference.} Original section 0FM6: sources/stacks-derham/pinned-original-source.tex, source lines 671--671. Navigation pointer for \textbackslash{}ref\{section-hodge-to-de-rham\} used inside the counted span (line 6348); the argument that consumes it is retained in the span.\par\medskip
\begin{proposition}
\label{proposition-relative-poincare-duality}
Let $S$ be a quasi-compact and quasi-separated scheme. Let $f : X \to S$
be a proper smooth morphism of schemes all of whose fibres are nonempty
and equidimensional of dimension $n$. There exists an
$\mathcal{O}_S$-module map
$$
t : R^{2n}f_*\Omega^\bullet_{X/S} \longrightarrow \mathcal{O}_S
$$
unique up to precomposing by multiplication by a unit of
$H^0(X, \mathcal{O}_X)$ with the following property: the pairing
$$
Rf_*\Omega^\bullet_{X/S}
\otimes_{\mathcal{O}_S}^\mathbf{L}
Rf_*\Omega^\bullet_{X/S}[2n]
\longrightarrow
\mathcal{O}_S, \quad
(\xi, \xi') \longmapsto t(\xi \cup \xi')
$$
is a perfect pairing of perfect complexes on $S$.
\end{proposition}

\begin{proof}
The proof is exactly the same as the proof of
Proposition \ref{proposition-poincare-duality}.

\medskip\noindent
By the relative Hodge-to-de Rham spectral sequence
$$
E_1^{p, q} = R^qf_*\Omega^p_{X/S} \Rightarrow R^{p + q}f_*\Omega^\bullet_{X/S}
$$
(Section \ref{section-hodge-to-de-rham}), the vanishing
of $\Omega^i_{X/S}$ for $i > n$, the vanishing in for example Limits, Lemma
\ref{limits-lemma-higher-direct-images-zero-above-dimension-fibre}
and the results of Lemmas \ref{lemma-relative-bottom-part-degenerates} and
\ref{lemma-relative-top-part-degenerates}
we see that $R^0f_*\Omega_{X/S} = R^0f_*\mathcal{O}_X$
and $R^nf_*\Omega^n_{X/S} = R^{2n}f_*\Omega^\bullet_{X/S}$.
More precisely, these identifications come from the maps
of complexes
$$
\Omega^\bullet_{X/S} \to \mathcal{O}_X[0]
\quad\text{and}\quad
\Omega^n_{X/S}[-n] \to \Omega^\bullet_{X/S}
$$
Let us choose $t : R^{2n}f_*\Omega_{X/S} \to \mathcal{O}_S$
which via this identification corresponds to a $t$ as in
Lemma \ref{lemma-relative-duality-hodge}.

\medskip\noindent
Let us abbreviate $\Omega^\bullet = \Omega^\bullet_{X/S}$.
Consider the map (\ref{equation-wedge}) which in our situation reads
$$
\wedge :
\text{Tot}(\Omega^\bullet \otimes_{f^{-1}\mathcal{O}_S} \Omega^\bullet)
\longrightarrow
\Omega^\bullet
$$
For every integer $p = 0, 1, \ldots, n$ this map annihilates the subcomplex
$\text{Tot}(\sigma_{> p} \Omega^\bullet \otimes_{f^{-1}\mathcal{O}_S}
\sigma_{\geq n - p} \Omega^\bullet)$ for degree reasons.
Hence we find that the restriction of $\wedge$ to the subcomplex
$\text{Tot}(\Omega^\bullet \otimes_{f^{-1}\mathcal{O}_S}
\geq_{n - p}\Omega^\bullet)$ factors through a map of complexes
$$
\gamma_p :
\text{Tot}(\sigma_{\leq p} \Omega^\bullet \otimes_{f^{-1}\mathcal{O}_S}
\sigma_{\geq n - p} \Omega^\bullet)
\longrightarrow
\Omega^\bullet
$$
Using the same procedure as in Section \ref{section-cup-product} we obtain
relative cup products
$$
Rf_*\sigma_{\leq p} \Omega^\bullet
\otimes_{\mathcal{O}_S}^\mathbf{L}
Rf_*\sigma_{\geq n - p}\Omega^\bullet
\longrightarrow
Rf_*\Omega^\bullet
$$
We will prove by induction on $p$ that these cup products via $t$
induce perfect pairings between $Rf_*\sigma_{\leq p} \Omega^\bullet$
and $Rf_*\sigma_{\geq n - p}\Omega^\bullet[2n]$. For $p = n$
this is the assertion of the proposition.

\medskip\noindent
The base case is $p = 0$. In this case we have
$$
Rf_*\sigma_{\leq p}\Omega^\bullet = Rf_*\mathcal{O}_X
\quad\text{and}\quad
Rf_*\sigma_{\geq n - p}\Omega^\bullet[2n] = Rf_*(\Omega^n[-n])[2n] =
Rf_*\Omega^n[n]
$$
In this case we simply obtain the pairing
between $Rf_*\mathcal{O}_X$ and $Rf_*\Omega^n[n]$ of
Lemma \ref{lemma-relative-duality-hodge} and the result is true.

\medskip\noindent
Induction step. Say we know the result is true for $p$. Then
we consider the distinguished triangle
$$
\Omega^{p + 1}[-p - 1] \to
\sigma_{\leq p + 1}\Omega^\bullet \to
\sigma_{\leq p}\Omega^\bullet \to
\Omega^{p + 1}[-p]
$$
and the distinguished triangle
$$
\sigma_{\geq n - p}\Omega^\bullet \to
\sigma_{\geq n - p - 1}\Omega^\bullet \to
\Omega^{n - p - 1}[-n + p + 1] \to
(\sigma_{\geq n - p}\Omega^\bullet)[1]
$$
Observe that both are distinguished triangles in the homotopy category
of complexes of sheaves of $f^{-1}\mathcal{O}_S$-modules; in particular the
maps $\sigma_{\leq p}\Omega^\bullet \to \Omega^{p + 1}[-p]$ and
$\Omega^{n - p - 1}[-d + p + 1] \to (\sigma_{\geq n - p}\Omega^\bullet)[1]$
are given by actual maps of complexes, namely using the differential
$\Omega^p \to \Omega^{p + 1}$ and the differential
$\Omega^{n - p - 1} \to \Omega^{n - p}$.
Consider the distinguished triangles associated gotten from these
distinguished triangles by applying $Rf_*$
$$
\xymatrix{
Rf_*\sigma_{\leq p}\Omega^\bullet \ar[d]_a \\
Rf_*\Omega^{p + 1}[-p - 1] \ar[d]_b \\
Rf_*\sigma_{\leq p + 1}\Omega^\bullet \ar[d]_c \\
Rf_*\sigma_{\leq p}\Omega^\bullet \ar[d]_d \\
Rf_*\Omega^{p + 1}[-p - 1]
}
\quad\quad
\xymatrix{
Rf_*\sigma_{\geq n - p}\Omega^\bullet \\
Rf_*\Omega^{n - p - 1}[-n + p + 1] \ar[u]_{a'} \\
Rf_*\sigma_{\geq n - p - 1}\Omega^\bullet \ar[u]_{b'} \\
Rf_*\sigma_{\geq n - p}\Omega^\bullet \ar[u]_{c'} \\
Rf_*\Omega^{n - p - 1}[-n + p + 1] \ar[u]_{d'}
}
$$
We will show below that the pairs $(a, a')$, $(b, b')$, $(c, c')$, and
$(d, d')$ are compatible with the given pairings. This means we obtain a
map from the distinguished triangle on the left to the distinguished triangle
obtained by applying $R\SheafHom(-, \mathcal{O}_S)$ to the distinguished
triangle on the right. By induction and Lemma \ref{lemma-duality-hodge}
we know that the pairings constructed above between the
complexes on the first, second, fourth, and fifth
rows are perfect, i.e., determine isomorphisms after taking duals.
By Derived Categories, Lemma \ref{derived-lemma-third-isomorphism-triangle}
we conclude the pairing between the complexes in the middle row
is perfect as desired.

\medskip\noindent
Let $e : K \to K'$ and $e' : M' \to M$ be maps of objects
of $D(\mathcal{O}_S)$ and let
$K \otimes_{\mathcal{O}_S}^\mathbf{L} M \to \mathcal{O}_S$ and
$K' \otimes_{\mathcal{O}_S}^\mathbf{L} M' \to \mathcal{O}_S$
be pairings. Then we say these pairings are compatible if the
diagram
$$
\xymatrix{
K' \otimes_{\mathcal{O}_S}^\mathbf{L} M' \ar[d] &
K \otimes_{\mathcal{O}_S}^\mathbf{L} M'
\ar[l]^{e \otimes 1} \ar[d]^{1 \otimes e'} \\
\mathcal{O}_S &
K \otimes_{\mathcal{O}_S}^\mathbf{L} M \ar[l]
}
$$
commutes. This indeed means that the diagram
$$
\xymatrix{
K \ar[r] \ar[d]_e & R\SheafHom(M, \mathcal{O}_S) \ar[d]^{R\SheafHom(e', -)} \\
K' \ar[r] & R\SheafHom(M', \mathcal{O}_S)
}
$$
commutes and hence is sufficient for our purposes.

\medskip\noindent
Let us prove this for the pair $(c, c')$. Here we observe simply
that we have a commutative diagram
$$
\xymatrix{
\text{Tot}(\sigma_{\leq p} \Omega^\bullet \otimes_{f^{-1}\mathcal{O}_S}
\sigma_{\geq n - p} \Omega^\bullet) \ar[d]_{\gamma_p} &
\text{Tot}(\sigma_{\leq p + 1} \Omega^\bullet \otimes_{f^{-1}\mathcal{O}_S}
\sigma_{\geq n - p} \Omega^\bullet) \ar[l] \ar[d] \\
\Omega^\bullet &
\text{Tot}(\sigma_{\leq p + 1} \Omega^\bullet \otimes_{f^{-1}\mathcal{O}_S}
\sigma_{\geq n - p - 1} \Omega^\bullet) \ar[l]_-{\gamma_{p + 1}}
}
$$
By functoriality of the cup product we obtain commutativity of the
desired diagram.

\medskip\noindent
Similarly for the pair $(b, b')$ we use the commutative diagram
$$
\xymatrix{
\text{Tot}(\sigma_{\leq p + 1} \Omega^\bullet \otimes_{f^{-1}\mathcal{O}_S}
\sigma_{\geq n - p - 1} \Omega^\bullet) \ar[d]_{\gamma_{p + 1}} &
\text{Tot}(\Omega^{p + 1}[-p - 1] \otimes_{f^{-1}\mathcal{O}_S}
\sigma_{\geq n - p - 1} \Omega^\bullet) \ar[l] \ar[d] \\
\Omega^\bullet &
\Omega^{p + 1}[-p - 1]
\otimes_{f^{-1}\mathcal{O}_S}
\Omega^{n - p - 1}[-n + p + 1] \ar[l]_-\wedge
}
$$

\medskip\noindent
For the pairs $(d, d')$ and $(a, a')$ we use the commutative diagram
$$
\xymatrix{
\Omega^{p + 1}[-p] \otimes_{f^{-1}\mathcal{O}_S}
\Omega^{n - p - 1}[-n + p] \ar[d] &
\text{Tot}(\sigma_{\leq p}\Omega^\bullet \otimes_{f^{-1}\mathcal{O}_S}
\Omega^{n - p - 1}[-n + p]) \ar[l] \ar[d] \\
\Omega^\bullet &
\text{Tot}(\sigma_{\leq p}\Omega^\bullet \otimes_{f^{-1}\mathcal{O}_S}
\sigma_{\geq n - p}\Omega^\bullet) \ar[l]
}
$$

\medskip\noindent
We omit the argument showing the uniqueness of $t$ up to
precomposing by multiplication by a unit in $H^0(X, \mathcal{O}_X)$.
\end{proof}

\end{document}
